\documentclass[a4paper, 11pt]{article}
\usepackage{cite}

\title{Paper Format Checker Test Document}
\author{Test Author}

\begin{document}
\maketitle

\section{Introduction}

% --- comma-space: 半角カンマの後にスペースなし（検出されるべき） ---
This sentence has a bad comma,without a space after it.
We tested cases A,B,C without spacing.

% comma-space: 正しい例（検出されないべき）
This sentence has a good comma, with a space after it.

% --- zenkaku-comma: 全角読点「、」の使用（検出されるべき） ---
本研究では、提案手法の有効性を示す、実験を行った。

% zenkaku-comma: 正しい例（全角カンマ「，」を使用、検出されないべき）
本研究では，提案手法の有効性を示す，実験を行った．

% --- hankaku-comma: 英字の後の全角ピリオド「．」（検出されるべき） ---
The accuracy was high．
The result was significant．

% hankaku-comma: 正しい例（半角ピリオド、検出されないべき）
The accuracy was high.

\section{Related Work}

% --- cite-error: ピリオドの後ろに \cite{}（検出されるべき） ---
This approach is widely used．\cite{smith2020}
A similar method was proposed.\cite{tanaka2019}

% cite-error: 正しい例（引用はピリオドの前、検出されないべき）
This approach is widely used~\cite{smith2020}．

\section{Conclusion}

% 複合的に問題のある行
我々は,提案手法を評価した、その結果は良好であった．\cite{yamada2021}

\bibliographystyle{plain}
\bibliography{references}

\end{document}
